% !TITLE 拟行路难·其四
% !AUTHOR 鲍照
% !DYNASTY 魏晋
% !GENRE 七言古诗
% !ORDER 8

\begin{poem}
泻\textsuperscript{1}水置平地，各自东西南北流。
人生亦有命，安能行叹复坐愁？
酌\textsuperscript{2}酒以自宽，举杯断绝\textsuperscript{3}歌路难。
心非木石岂无感？吞声踯躅\textsuperscript{4}不敢言。
\end{poem}

\begin{annotations}
\notetext{1}{泻：泼洒，倾倒。}
\notetext{2}{酌：酌（zhuó），倒酒，斟酒。}
\notetext{3}{断绝：断开、止息。这里指歌声戛然而止，唱不下去。}
\notetext{4}{踯躅：踯躅（zhí zhú），徘徊，走来走去。形容内心情感剧烈冲突、难以决断时的身体动作。}
\end{annotations}

\begin{appreciation}
《拟行路难》是鲍照的代表作，一共十八首，这是第四首。它写的是一口气，咽下去又吐不出来的那种，是中国诗里写"憋屈"最透的诗篇之一。

泻水置平地，各自东西南北流——把水倒在平地上，水往东西南北自己流。这个比喻妙极了：水往哪流，取决于地势，不取决于水。人也一样：人往哪里走，不全看才华，还看出身。鲍照的时代叫门阀制度，世家把持官位，寒门子弟再有本事也上不去。鲍照自己就是寒门出身。

人生亦有命，安能行叹复坐愁——人生自有命运，何必走路也叹、坐下也愁？话说得洒脱，其实是自我安慰，而且安慰不动——真认命，就不会写这首诗了。酌酒以自宽，举杯断绝歌路难——倒酒宽慰自己，举起杯，行路难的歌声戛然而止。不是不想唱，是越唱越堵，唱不下去了。

心非木石岂无感？吞声踯躅不敢言——心不是木头石头，怎么会没有感触？只是把声音咽回去，原地徘徊，不敢说。吞声是把话咽回肚子里，踯躅是来来回回地走。这是全诗最痛的地方：他有不平，有愤怒，却不敢言。

屈原在《天问》里敢把问题抛给上天，一口气问三百多个；鲍照连话都不敢说。不是他胆小，是时代不同：屈原觉得自己有资格问，鲍照连问的资格都没有。这个题目后来李白也写过：长风破浪会有时，直挂云帆济沧海。鲍照咽下去的话，李白替他吼了出来。

把这首诗选进来，是想让你知道：这世上曾经有人连话都不敢说。所以你别学"吞声"，遇事有不平，好好说出来，就比鲍照赢了。
\end{appreciation}
